\emph{Setting.} By Fact \claimref{pairing.imported-schwarzschild-connection-curvature} the static gauge has $\hb^{-1}(e_{\mu}) = g_{\mu}$ with $g_{t} = \sqrt f\,\gamma_{0}$, $g_{r} = \gamma_{1}/\sqrt f$, $g_{\vartheta} = r\gamma_{2}$, $g_{\varphi} = r\sin\vartheta\,\gamma_{3}$ and the rotation-gauge field of the Fact; the field is torsion-free, $D_{\mu}g_{\nu} = D_{\nu}g_{\mu}$ with $D_{\mu}M = \partial_{\mu}M + \Omega(e_{\mu}) \times M$ (\claimref{obj.synthesis.gtg-flat-spacetime-gauge-fields}), and its Ricci vector vanishes (both confirmed by the check at its points). Derived metric components are $g_{\mu} \cdot g_{\nu}$.

\emph{First-order fields.} Put $g_{\mu} \mapsto g_{\mu} + \epsilon k_{\mu}$ and $\Omega(e_{\mu}) \mapsto \Omega(e_{\mu}) + \epsilon\Omega_{1}(e_{\mu})$ with $\epsilon^{2} = 0$. The torsion-free condition at order $\epsilon$ is
\[
\Omega_{1}(e_{\mu}) \cdot g_{\nu} - \Omega_{1}(e_{\nu}) \cdot g_{\mu} = -\bigl(D_{\mu}k_{\nu} - D_{\nu}k_{\mu}\bigr),
\]
a linear system for the $24$ components of $\Omega_{1}$ whose matrix is that of the background torsion map, which is invertible (it is the map that determines $\Omega$ from the coframe). So $\Omega_{1}$ is unique. The field strength $R(e_{\mu} \wedge e_{\nu}) = \partial_{\mu}\Omega(e_{\nu}) - \partial_{\nu}\Omega(e_{\mu}) + \Omega(e_{\mu}) \times \Omega(e_{\nu})$ has first-order part $R_{1}(e_{\mu} \wedge e_{\nu}) = D_{\mu}\Omega_{1}(e_{\nu}) - D_{\nu}\Omega_{1}(e_{\mu})$, and the covariant curvature on frame planes, $\mathcal{R}(\gamma_{a} \wedge \gamma_{b}) = -R(\hb(\gamma_{a}) \wedge \hb(\gamma_{b}))$ in the store's sign, has first-order part $-R_{1}(\hb(\gamma_{a}) \wedge \hb(\gamma_{b})) - R(\hb_{1}(\gamma_{a}) \wedge \hb(\gamma_{b})) - R(\hb(\gamma_{a}) \wedge \hb_{1}(\gamma_{b}))$ with $\hb_{1}$ the first-order change of $\hb$ (minus the background inverse coframe applied on both sides of $k$). The first-order Ricci vector is its contraction, $\mathcal{R}_{1}(\gamma_{b}) = \sum_{a}\gamma^{a} \cdot \mathcal{R}_{1}(\gamma_{a} \wedge \gamma_{b})$. Because the background Ricci vector vanishes, $\mathcal{R}_{1}(\gamma_{b})$ is unchanged by a first-order local rotation of the frame, so it depends on $k$ only through the first-order metric $g_{\mu} \cdot k_{\nu} + k_{\mu} \cdot g_{\nu}$; the choice of $k$ realising a given metric change is immaterial.

\emph{(1).} For the stated $k_{\varphi}$ the first-order metric is $2\epsilon\,S\,(h_{0}\,dt\,d\varphi + h_{1}\,dr\,d\varphi)$ (using $\gamma_{0}^{2} = 1$, $\gamma_{1}^{2} = -1$), the axisymmetric odd-parity field in the Regge-Wheeler gauge with angular function $S = \sin\vartheta\,P'$, which satisfies $S'' = \cot\vartheta\,S' - LS$ by the equation of $P$. The check carries out the computation of the previous paragraph exactly: it solves the linear system for $\Omega_{1}$ with jets, forms $\mathcal{R}_{1}$, and compares every component of every $\mathcal{R}_{1}(\gamma_{b})$ with the displayed expressions. Both sides are linear in the second-order jets of $h_{0}$ and $h_{1}$ at the point and in the values $(S, S')$ there ($S''$ being eliminated by $S'' = \cot\vartheta\,S' - LS$), so it suffices to compare on a basis: the six monomials $t^{p}(r - r_{0})^{q}$, $p + q \leq 2$, for $h_{0}$ with $h_{1} = 0$ and for $h_{1}$ with $h_{0} = 0$, and the data $(S, S') = (1, 0)$ and $(0, 1)$. For each basis element the comparison is run with the jet coefficients in the field of rational functions of $z = \sqrt f$, where $r = 2M/(1 - z^{2})$ and $\partial_{r}z = f'/(2z)$ is rational in $z$; it is then an identity of rational functions of $z$ and holds at every radius $r > 2M$. This is done at the three polar angles with $(\sin\vartheta, \cos\vartheta) = (3/5, 4/5), (5/13, 12/13), (15/17, -8/17)$, whose cotangents $4/3$, $12/5$, $-8/15$ are distinct, and for $L = 7/2$ and $L = 6$ at each angle: $144$ cases, all in agreement.

\emph{Why the grid decides the identity.} Dependence on $L$: $L$ enters only through the elimination of $S''$, once, so both sides are of degree at most one in $L$, and two values of $L$ at each angle decide every $L$. Dependence on the polar angle (the degree bound; argued): treat $s = \sin\vartheta$ and $c = \cos\vartheta$ as formal functions with $s' = c$, $c' = -s$. The background coframe involves the angle only through $g_{\varphi} = rs\,\gamma_{3}$, and the substitution $(s, c, S) \mapsto (as, ac, aS)$ with $a$ constant is the rescaling $\varphi \mapsto a\varphi$ of the background label, under which the frame components of $\mathcal{R}_{1}$ do not change. So for fixed radius and each basis jet the coefficients of $S$ and of $S'$ in each frame component are homogeneous of degree $-1$ in $(s, c)$. $\mathcal{R}_{1}$ is of total derivative order two, and each derivative of $s$ or of the angular data brings at most one factor $c$ (the elimination of $S''$ brings one factor $c/s$ with $S'$), so the degree in $c$ is at most two. Hence each coefficient lies in the span of $1/s$, $c/s^{2}$, $c^{2}/s^{3}$, that is $s$ times it is a polynomial of degree at most two in $\cot\vartheta$, and the same holds for the displayed right-hand sides ($S/\sin\vartheta$ and $(S' - 2\cot\vartheta\,S)/\sin\vartheta$). Three angles with distinct cotangents therefore decide the identity at every polar angle. Dependence on $M$: the equations are homogeneous under $(t, r, M) \mapsto (ct, cr, cM)$ with $h_{0}, h_{1}$ fixed, so $M = 1$ suffices. The degree bound is an argument, not a computation; as controls the check also runs the comparison at two rational points with generic jets and $L = 6, 12, 7/2$, and as an identity of rational functions of $\tan(\vartheta/2)$ (every polar angle) at two fixed radii with generic jets, where it holds. The expressions are those of Regge and Wheeler (from memory), an independent control.

\emph{(2).} From the definitions, $\partial_{t}h_{0} = f\partial_{r}(fh_{1}) + fE_{\vartheta}$. Substituting in $E_{r}$ and writing $fh_{1} = rQ$, so that $\partial_{r}(fh_{1}) = Q + rQ'$:
\[
E_{r} = \frac{r}{f}\partial_{t}^{2}Q - \partial_{r}\bigl(fQ + rfQ'\bigr) - \partial_{r}(fE_{\vartheta}) + \frac{2f}{r}\bigl(Q + rQ'\bigr) + \frac{2f}{r}E_{\vartheta} + \frac{(L-2)}{r}Q .
\]
Now $\partial_{r}(fQ + rfQ') = f'Q + 2fQ' + r(fQ')'$, so the $Q'$ terms cancel and, multiplying by $f/r$,
\[
\frac{f}{r}\Bigl(E_{r} + \partial_{r}(fE_{\vartheta}) - \frac{2f}{r}E_{\vartheta}\Bigr) = \partial_{t}^{2}Q - f(fQ')' + f\Bigl(\frac{L - 2}{r^{2}} + \frac{2f}{r^{2}} - \frac{f'}{r}\Bigr)Q ,
\]
and with $f = 1 - 2M/r$, $f' = 2M/r^{2}$ the bracket is $L/r^{2} - 6M/r^{3}$; $f(fQ')' = \partial_{r_{*}}^{2}Q$. For the second identity, $\partial_{t}E_{t} + f\partial_{r}E_{r}$ has no third derivatives; subtracting $-(2f/r)E_{r}$ leaves $\partial_{t}h_{0}\bigl(-L/r^{2} + 4M/r^{3} + 2f/r^{2}\bigr) + (L-2)f\,\partial_{r}(fh_{1})/r^{2} = -(L-2)(f/r^{2})E_{\vartheta}$. Both identities are verified by the check for generic jets. Hence $E_{r} = E_{\vartheta} = 0$ imply the Regge-Wheeler equation. Conversely let $Q$ solve it on a channel with $\partial_{t}^{2} = -\omega^{2}$, $\omega \neq 0$; put $h_{1} = rQ/f$ and $h_{0} = -\omega^{-2}\partial_{t}\bigl(f\partial_{r}(rQ)\bigr)$, the solution of $\partial_{t}h_{0} = f\partial_{r}(rQ)$ in the channel. Then $E_{\vartheta} = 0$, the first identity gives $E_{r} = (r/f)\bigl(\partial_{t}^{2}Q - \partial_{r_{*}}^{2}Q + V_{\mathrm{o}}Q\bigr) = 0$, and the second gives $\partial_{t}E_{t} = 0$; since $E_{t}$ lies in the channel, $-\omega^{2}E_{t} = \partial_{t}^{2}E_{t} = 0$ and $E_{t} = 0$. By (1) all of $\mathcal{R}_{1}(\gamma_{b})$ vanish. The check runs one reconstruction with a polynomial $Q$.

\emph{(3).} The stated $k_{\mu}$ give the first-order metric $-fH_{0}P\,dt^{2} - 2H_{1}P\,dt\,dr - f^{-1}H_{2}P\,dr^{2} - r^{2}KP\,(d\vartheta^{2} + \sin^{2}\vartheta\,d\varphi^{2})$, the even-parity field in the Regge-Wheeler gauge with the signs of the mostly-minus signature. The check computes $\mathcal{R}_{1}(\gamma_{b})$ and its first $t$- and $r$-derivatives as linear functionals on the forty coefficients of the third-order jets of $(H_{0}, H_{1}, H_{2}, K)$ in $(t,r)$, with $P$ the jet of the solution of its equation with generic initial values, and finds rank $18$; adjoining the functional $\partial_{t}^{2}\Psi - \partial_{r_{*}}^{2}\Psi + V_{\mathrm{e}}\Psi$ leaves the rank $18$, and adjoining the same expression with $\lambda + 1$ in place of $\lambda$ in $V_{\mathrm{e}}$ raises it to $19$. So at those points the Zerilli expression is a combination of the vacuum equations and their first derivatives. The coefficients of the combination are not exhibited, and the statement is sampled in the point and in $L$ (the harmonic value $6$ and the generic value $7/2$). The function $\Psi$ and the potential are those of Zerilli and Moncrief (from memory). The converse is (C3).

\emph{(4).} With $D := \partial_{r_{*}} = f\partial_{r}$: $(D + W)(D - W) = D^{2} - (DW) - W^{2}$ and $(D - W)(D + W) = D^{2} + (DW) - W^{2}$ as operators. The two identities for the potentials are identities of rational functions of $r$ (verified exactly for three values of $M$ and $L = 6, 12, 20, 7/2$; both sides are rational in $M$, $\lambda$, $r$). They give $-D^{2} + V_{\mathrm{o}} = -(D - W)(D + W) - c^{2}$ and $-D^{2} + V_{\mathrm{e}} = -(D + W)(D - W) - c^{2}$, so $(D + W)(-D^{2} + V_{\mathrm{o}}) = (-D^{2} + V_{\mathrm{e}})(D + W)$, and $\partial_{t}$ commutes with $D$ and $W$. If $Q$ solves the odd equation with $\partial_{t}^{2}Q = -\omega^{2}Q$, then $(-D^{2} + V_{\mathrm{o}})Q = \omega^{2}Q$ and $(D - W)(D + W)Q = -(\omega^{2} + c^{2})Q$, which gives the inverse; $(D + W)Q = 0$ forces $(-D^{2} + V_{\mathrm{o}})Q = -c^{2}Q$, that is $\partial_{t}^{2}Q = c^{2}Q$. At $r = 2M$, $f = 0$ and $W = c$.

\emph{(5).} $V_{\mathrm{o}}/f = L/r^{2} - 6M/r^{3}$ is $(L - 3)/(4M^{2})$ at $r = 2M$. For $V_{\mathrm{e}}/f$ the numerator at $r = 2M$ is $M^{3}(16\lambda^{3} + 40\lambda^{2} + 36\lambda + 18) = 2M^{3}(2\lambda + 3)(4\lambda^{2} + 4\lambda + 3)$ and the denominator $8M^{5}(2\lambda + 3)^{2}$; the denominator $r^{3}(\lambda r + 3M)^{2}$ has no zero on $r > 0$ for $\lambda \geq 0$. By clause (4) of Theorem \claimref{synthesis.horizon-channels-of-scalar-dirac-and-maxwell-fields-are-two-dimensional}, $f = C(r)e^{2\kappa r_{*}}$ with $C$ regular and positive at the horizon, so $V_{\mathrm{o}}, V_{\mathrm{e}} = O(e^{2\kappa r_{*}})$ with nonzero leading coefficient, and each master equation tends to $\partial_{t}^{2}u - \partial_{r_{*}}^{2}u = 0$, solved by functions of $t - r_{*}$ and of $t + r_{*}$. The count of two channels per $(\omega, \ell, m)$ is the two parities, by (C1). The approach is not uniform in $\ell$ because $V/f$ grows as $L$. That every linearised field is covered uses (C1) and (C2); for $\ell = 0, 1$ the Regge-Wheeler gauge and the master functions are not available ($L - 2 = 0$ at $\ell = 1$ removes $h_{1}$ from $E_{r}$'s algebraic term, and the perturbations are changes of mass and spin and gauge), so the statement is for $\ell \geq 2$.
